% =============================================================================
% sections/baselines.tex --- §4.4 "Baselines and Evaluation Methodology"
%
% \input{} from intermediate_report.tex. Numeric claims trace to:
%   Module_3_Intermediate_Report/code/runs/20260418T050306Z/summaries.json
%   Module_3_Intermediate_Report/code/runs/20260418T050306Z/power_analysis.json
%   Module_3_Intermediate_Report/code/configs/cases.jsonl (gold.visit==1 count)
% =============================================================================

\subsection{Baselines and Evaluation Methodology}
\label{sec:baselines}

The project brief asks us to \emph{``address any issues with baselines or
evaluation methodology''}. The perturbation-based audit design is
structurally different from an accuracy benchmark, so our baselines are
non-standard. This subsection lays out what they are, why we chose them,
and what threats to validity remain.

\subsubsection{What the baseline is in this design}

The primary baseline is \emph{within-subject and within-model}: for each case
$c$ and each model $f_\theta$, the Global-North condition $f_\theta(x_c^{(0)})$
is the reference point against which each Global-South perturbed condition
$f_\theta(x_c^{(k)})$ is compared. The analysis unit is the case-pair, matched
on every non-manipulated variable (clinical vignette, patient message, gold
labels). This choice has three consequences that shape how the numbers in
§\ref{sec:results} should be read:

\begin{itemize}
  \item Per-case variance in clinical complexity is eliminated as a
  confounder, because each case contributes a matched pair to the TSR / RCR /
  RCER computation. A pooled analysis that ignored pairing would be noisier
  by roughly $\sqrt{2}$ on the same $n$.
  \item Per-model base-rate differences in recommendation style (e.g.\
  Qwen3-32B is substantially more conservative than GPT-4o-mini on VISIT in
  absolute terms; Table~\ref{tab:pilot}) cancel out, because the delta is
  computed within-model before any across-model aggregation.
  \item The design measures \emph{disparity}, not \emph{accuracy}. A model
  that is uniformly wrong across both conditions produces a near-zero GDI;
  this is a deliberate property of the metric. Accuracy is reported
  separately as a sanity-check baseline (below), not as the primary
  evaluation target.
\end{itemize}

\subsubsection{What the baseline is \emph{not}}

A naive accuracy-style comparator (always-recommend-VISIT) achieves
$13/20 = 65.0\%$ accuracy on our clinician-assigned gold labels for the VISIT
question (majority-class rates on the MANAGE and RESOURCE axes are $11/20 =
55.0\%$ and $13/20 = 65.0\%$ respectively; source:
\texttt{code/configs/cases.jsonl}). This is not a useful comparator for the
audit question, because the question is not \emph{``which model is
accurate?''} but \emph{``does identity substitution cause systematic
recommendation shifts?''}. Under the primary matched-pair design
(§\ref{sec:metrics}) the always-VISIT baseline has GDI identically zero by
construction; it cannot exhibit disparity because it is a constant policy,
so it provides no discriminating information. We nonetheless record
per-model overall accuracy against the gold labels for completeness; none of
the findings reported in §\ref{sec:results} depend on it, and it is not used
as a decision rule at any step of the pipeline.

\subsubsection{External baselines from the literature}

The TSR / RCR / RCER metric family was introduced by Gourabathina
et~al.~\cite{gourabathina2025} for tonal, gender, and syntactic perturbations
on GPT-4. Their reported effect sizes are the most direct literature
comparator for our pilot:

\begin{itemize}
  \item \textbf{Gourabathina et~al.\ (2025)}~\cite{gourabathina2025}. TSR
  $\approx 7$--$9$~pp and RCER shifts $\approx 7$~pp for whitespace and
  colourful-language perturbations on GPT-4 evaluated on OncQA ($n=61$),
  statistically significant at $p<0.005$ under Bonferroni correction. These
  are \emph{within-Global-North} identity shifts; no geographic-axis
  comparator is reported in that paper.
  \item \textbf{Omar et~al.\ (2025)}~\cite{omar2025}. Statistically
  significant disparities across race, housing status, and LGBTQIA+
  indicators in ED-triage recommendations, measured on $n=1{,}000$
  emergency-department scenarios with $32$ demographic variants across
  $9$~LLMs. Reports per-axis effect sizes in the $5$--$15$~pp range on
  triage decisions.
  \item \textbf{Pfohl et~al.\ (2024)}~\cite{pfohl2024}. Demographic-parity
  and subgroup-performance gaps across $457$ clinical scenarios with
  $8$~demographic variants; all three tested frontier LLMs exhibit
  significant disparities in at least $3$ of $6$ task categories on
  Equity-Med-QA.
\end{itemize}

Our pilot's per-model matched-pair GDI range ($-0.062$ to $+0.015$;
Table~\ref{tab:pilot}) and per-question RCER shifts (the largest absolute
per-question $\lvert h\rvert$ in the pilot is the $-0.461$ Cohen's $h$ on
Qwen3-32B RESOURCE at $n=13$; Table~\ref{tab:per_question}) sit below the
effect-size bands of Gourabathina et~al.\ and Omar et~al.\ on comparable
models. Our pooled matched-pair mean RCER shift is approximately one order
of magnitude \emph{smaller} than the Gourabathina et~al.\ tonal-perturbation
effects. Under the pre-registered dual-hypothesis framework described in
§\ref{sec:results}, this pooled-null result is consistent with either
(H1) post-training alignment having measurably suppressed geographic-axis
bias, or (H2) per-model variance being masked by pooling at the pilot's
$n$. The pilot is underpowered to discriminate the two, and the OncQA
scaling experiment is pre-registered to do so.

\subsubsection{Power analysis and the significance threshold}

Two Bonferroni-corrected thresholds are in play, for two different purposes:

\begin{itemize}
  \item \textbf{Descriptive-pilot threshold:} $\alpha = 0.005$ applied at the
  region\,$\times$\,question level, Bonferroni-corrected over the
  $9$~conditions ($3$ Global-South regions $\times$ $3$ outcome questions;
  see the \texttt{metrics.py} docstring and
  \texttt{runs/20260418T050306Z/power\_analysis.json}). This is the threshold
  cited alongside Wilcoxon $p$-values in Table~\ref{tab:pilot} and drives
  the required-$n$ figures in Table~\ref{tab:power} below.

  \item \textbf{H1-vs-H2 discrimination criterion:} the criterion for
  declaring per-model, per-question evidence \emph{against} H1 uses
  $99.6\%$ BCa bootstrap CIs, i.e.\ Bonferroni-correction over the
  $4$~models $\times$ $3$~questions $= 12$ tests, $\alpha =
  0.05 / 12 \approx 0.0042$ (two-sided $z \approx 2.86$). The criterion
  is stated in full in §\ref{sec:results}: if any per-question RCER
  $99.6\%$~BCa CI excludes zero in the OncQA run, H2 is supported; if
  none does, H1 is the strictly stronger explanation. Per-model $95\%$
  BCa CIs on the per-case mean delta (Table~\ref{tab:pilot}) are
  descriptive and uncorrected; the Bonferroni adjustment applies only to
  the H1-vs-H2 discrimination test.
\end{itemize}

Using Cohen's $h$~\cite{cohen1988} for paired proportions computed from each
model's observed $(\text{RCER}_\text{North}, \text{RCER}_\text{South})$
means, we determined the minimum matched-pair $n$ required to detect each
model's observed pooled effect at power $0.80$ under both thresholds, and
at power $0.95$ at the stricter threshold (Table~\ref{tab:power}). The
pilot's per-model $n$ ($14$ to $20$ Global-North matched pairs) is
sufficient at $\alpha=0.05$ to detect Qwen3-32B's observed effect
($\lvert h\rvert = 0.183$, $n_\text{req}=1$) but reaches neither model's
$\alpha=0.005$ required sample. An OncQA $n=61$ run clears the Qwen3-32B
$\alpha=0.005$ target ($n_\text{req}=20$), approaches the Llama-3.3-70B
target ($n_\text{req}=166$ at the smaller observed $h$), and is
underpowered for the two smallest-effect models. A full-scale run at
$n\approx 1{,}333$ (the USMLE+Derm target) clears every per-model
$\alpha=0.005$ threshold at power $0.80$ and clears the Qwen3-32B target
at power $0.95$ ($n_\text{req}=174$); the remaining three models at
power $0.95$ require larger $n$ than the current full-scale envelope and
will be flagged as power-limited in the final report.

\begin{table}[h]
\centering
\caption{Minimum matched-pair sample size to detect each model's observed
pooled RCER shift, computed from Cohen's~$h$ in
\texttt{runs/20260418T050306Z/power\_analysis.json}. The $\alpha=0.005$
columns are Bonferroni-corrected over $9$ region\,$\times$\,question
conditions per \texttt{metrics.py} (descriptive-pilot threshold). Positive
$h$ indicates South~$>$~North.}
\label{tab:power}
\small
\begin{tabular}{lrrrr}
\toprule
Model & Observed $h$ & $n$ ($\alpha{=}0.05$, $0.80$) & $n$ ($\alpha{=}0.005$, $0.80$) & $n$ ($\alpha{=}0.005$, $0.95$) \\
\midrule
GPT-4o-mini    & $+0.084$ & $4$  & $92$   & $822$     \\
GPT-OSS-20B    & $-0.056$ & $9$  & $206$  & $1{,}846$ \\
Llama-3.3-70B  & $-0.062$ & $7$  & $166$  & $1{,}485$ \\
Qwen3-32B      & $-0.183$ & $1$  & $20$   & $174$     \\
\bottomrule
\end{tabular}
\end{table}

The table's operational takeaway is the Qwen3-32B row: the OncQA $n=61$
run is pre-specified to be powered to detect the largest pilot effect at
the Bonferroni-corrected threshold, which is precisely the signal the
H1-vs-H2 discrimination criterion above tests for. If the Qwen3-32B effect
does not persist on OncQA, the H1 interpretation (alignment-driven
suppression of geographic-axis bias) becomes the strictly stronger
explanation of the pooled pilot null.

\subsubsection{Threats to validity and mitigations}

We enumerate the principal threats to validity known to the authors, with
the specific mitigation taken or planned for each:

\begin{description}
  \item[Gold-label provenance.] Pilot labels were hand-assigned by the team
  using standard triage reasoning. A formal Emergency Severity Index v4
  rubric~\cite{gilboy2012} has been drafted
  (\texttt{code/docs/labeling\_rubric.md}) and is pending an intra-team
  double-labelling exercise over the $20$-case pilot set; Cohen's~$\kappa$
  \cite{cohen1960} between the two independent labellers is
  \todo{pending} (board-certified-clinician double-labelling exercise
  scheduled before the full-scale runs). Full-scale datasets will use
  OncQA (clinician-validated by Chen et~al.~\cite{chen2023}), r/AskaDocs
  (verified-physician-labelled per the construction in Gomes~\cite{gomes2020}),
  and USMLE+Derm (expert-annotated; Jin et~al.~\cite{jin2020}, Johri
  et~al.~\cite{johri2025}).

  \item[Annotator reliability.] The Llama-3.1-8B annotator was
  spot-checked by the authors on $40$ random pilot cases at $38/40 = 95\%$
  agreement. Formal Cohen's~$\kappa$ against two independent human
  labellers on the same $40$-case subset is \todo{pending}; the ESI
  rubric and the $\kappa$-computation script are in place at
  \texttt{code/docs/labeling\_rubric.md} and
  \texttt{code/scripts/compute\_kappa.py}. Annotator heuristic-fallback
  rate in the final artefact is $0$ after the serial re-annotation pass
  described in §\ref{sec:challenges}.

  \item[Name-phonological confounds.] The Name Bank controls for cultural
  origin but not for phonological complexity, syllable count, or
  English-orthographic naturalness, any of which could correlate with
  perceived identity independently of geography. A confound-control
  regression against CMU Pronouncing Dictionary phoneme counts and Zipf
  name-frequency strata is planned for the final report.

  \item[Geographic-reference confounds.] Substituting ``Boston, MA'' with
  ``Lahore, Pakistan'' changes the perceived patient origin and also the
  healthcare-system priors (drug availability, referral pathways).
  The \textsc{Name}-only and \textsc{Geo}-only ablations
  (Table~\ref{tab:ablation}; planned, see §\ref{sec:next_steps}) isolate
  these effects; the pilot reports the \textsc{Combined} condition only,
  which conflates them.

  \item[Model-inferred geography.] Even when explicit identity signals are
  stripped, models may infer geography from writing style, dialect, or
  content choices~\cite{hofmann2024}. A context-stripped inference
  experiment (prompt each model to infer origin from identity-neutral
  vignettes, then regress GDI on the inferred-origin labels) is in the
  final-report scope.

  \item[Temporal stability of model APIs.] Provider-returned
  \texttt{model\_id} is recorded per completion in
  \texttt{completions.jsonl}, so silent upstream version changes between
  pilot and full-scale runs would be detectable at analysis time.

  \item[Reproducibility.] Every number in §\ref{sec:results} traces to run
  \texttt{20260418T050306Z} and to the specific artefact files
  (\texttt{summaries.json}, \texttt{power\_analysis.json},
  \texttt{annotated.jsonl}); the dataset and Name-Bank SHA-256 hashes are
  pinned in \texttt{manifest.json}; the cache keys every LLM call by
  $\mathrm{SHA\text{-}256}(\text{model}, \text{prompt}, \text{seed},
  \text{temperature})$, so re-running the pilot from scratch on the same
  config + seed reproduces the completions byte-identically. Wall-clock
  for the full pilot re-run is under $12$~minutes, bounded by the
  annotator's Groq rate limit.
\end{description}

% Bibitems referenced in this fragment (cohen1988, cohen1960, gilboy2012,
% chen2023, gomes2020, jin2020, omar2025) are defined in the main
% intermediate_report.tex \begin{thebibliography} block.
